%Chapter XXIV   General Theory of Thin Plates and Shells

\section{Formulae relating to the curvature of surfaces}
For the investigations in the last Chapter the elements of the theory of
the curvature of surfaces are adequate. For the purpose of developing a
more general method of treatment of the problem of curved plates or shells
we shall require some further results of this theory. It seems best to
begin by obtaining these results.

Let \((\alpha ,\beta) \)
denote any two parameters by means of which the position of
a point on a surface can be expressed, so that the equations
\(\alpha=\) const.,
\(\beta=\) const., represent families of curves traced on the surface. Let
\(\chi \) be the angle between the tangents of these curves at any points.
\(\chi \) is in general
a function of
\((\alpha ,\beta) \).

\todo[inline]{footnote}
\[ x : R^2 \mapsto R^3,  ~~ < h, Dx' Dx h> \]
\[ (A,B) = (\norm{x,a},\norm{x,b}), A B \cos \chi =   \langle x,a , x,b \rangle \]
\[
Dx' Dx  =
\begin{bmatrix}
A^2  &  AB \cos \chi \\
AB \cos \chi & B^2
\end{bmatrix} =
\begin{bmatrix}
E & F \\
F & G
\end{bmatrix}.
\]

\todo[inline]{Transcribe again this time with MathPix}

The linear element \(ds\) of any curve traced on the surface is given by the formula
\begin{equation}\label{eq1}
 {(ds)}^2 = A^2 {(d \alpha)}^2   + B^2 {(d \beta)}^2 + 2AB \cos (\chi) ~~ d \alpha ~  d \beta
\end{equation}
where \(A\) and \(B\) are, in general, functions of
\((\alpha ,\beta)\). Let a right-handed system
of moving axes of \( (x, y, z) \)
be constructed so that the origin is at a point

 (0, 0)


of the surface, the axis of  is the normal to the surface at the origin, drawn
in a chosen sense, the axis of
x is the tangent to the curve
0 = const, which
passes through the origin, drawn in the sense of increase of
a, and the axis of
y is tangential to the surface, and at right angles to the axis of
x.

\todo[inline]{footnote}
footnote When the level set curves of \(\alpha \) and \(\beta \),
cut at right angles we suppose that the parameters
\((\alpha ,\beta) \)  and the positive sense of the normal to the surface, are
so chosen that the directions in which
\((\alpha ,\beta) \)  increase and this normal are the directions of a
right-handed system of axes.


When the origin of this triad of axes moves over the surface the directions of
the axes change. If
t represents the time, the components of velocity of the origin are
\[
A da/dt (1,0,0) +  B db/dt ( \cos \chi, \sin \chi, 0 )
\]
parallel to the instantaneous positions of the axes of
x, y, z. The components
af the angular velocity of the system of axes, referred to these same directions,
can be expressed in the forms

\[
(v1, v2, v3) = (p1,q1,r1) da/dt + (p2,q2,r2) db/dt
\]

% Pi da , d0 n da d0

in which the quantities p \dots are functions of o and 0.

\[
Dx ~~ [d \alpha; d \beta] ~~dt +  \cross{v}{x} +
\]

The quantities pi \ldots are connected with each other and with A, B,  by
the systems of equations (2) and (3) below. These results may be obtained
as follows:

 Let
x y z
denote the coordinates of a fixed point referred to the moving axes. Then
XyyZ
are functions of
\((\alpha ,\beta) \), and the conditions that the point remains fixed while the
axes move are the three equations

%dx da , dx dfi ( da  dff\ f da  d\ , a da j.d
%da dt 0/3
%da dt dp dt
%lt;i/3
%da dt  0/3
%/ da dp\ /da dp\  dgt;P .
%)+' ("o +'’ n) di
%dp f da dp\ / da dp\

Since these hold for all values of dajdt and dpjdty we have the six equatjons
% dx
%dx -A+riy-qiZ,
% da
%dy
% -cosx+rgy-a, -5sinx+/32“2,
% dz 0a
%Vi \

 The conditions of compatibility of these equations are three equations of the form
 () and ()  forming the differential coefficients, we may use the above

expressions for dx/dy \ldots The results must hold for all values of Xy z.

 The process just sketched leads to the equations
%3Pi 3P2 38 -

and

% ay?  da da 2 da ’oriP2 rjpj
% 1 fdA Bsinx 1 /dB (dA dB\ ' daj dA'
% J.sin5fV9o | + |8inx = |co8x.  (2) .(3)
To express the curvature of the surface we form the equations of the
normal at

%( + So, + 3/9) referred to the axes of
%x, y, z at
%(a, /9). The
direction-cosines of the normal are, with sufficient approximation,

%(gfjSa+jjS), - (p,3a + pa3/8), 1,
and the equations k,

%x - (ABa + BB0 cos y)  y - BBfi sin x
%{qM + qtB0) “-(j9,8+/),8/8)“’

It follows that the lines of curvature are given by the differential equation

%Pi (lt;) + 5 (jgt;8 cos X + sin x) (dY + |-4p + JS ( cos x +?i sin x)} dad 0,

(4) Ths Mts of equations (8) and (8) were obtained by
U. Codazzi, Paris, Him par di AOzantA, t. 27 (ISM)

and that the principal radii of curvature are the roots of the equation if

%(pi? -Ml) ~ i ~  iPi cos X + sin )) + ilB sin x  0. . . .(5)

From these results the equation of the indicatrix of the surface is easily
found to be

%"S " A +  = const ( 6 )

The measure of curvature is given by (5) and the third of (2) in the form

ABsmX\P

